\begin{tikzpicture}[
    node distance=1.2cm and 1.2cm,
    every node/.style={
        rectangle, draw, align=center, rounded corners=1pt,
        minimum height=1.2cm, font=\small, inner sep=5pt
    },
    every path/.style={-{Stealth[length=6pt]}, thick},
    efecto/.style={text width=4.6cm},
    problema/.style={text width=11.5cm, line width=0.9pt},
    causa/.style={text width=3.6cm},
    sub/.style={text width=2.5cm, font=\footnotesize},
    banda/.style={draw=none, font=\bfseries, anchor=east, align=right}
]

% ---------- Problema central ----------
\node (problema) [problema] {
El operario destina 2~h~43~min/mes al procedimiento manual\\
de dosificación de ingredientes secos, en el cual se acumula\\
una diferencia positiva de 1,976~kg de harina, 336~g de azúcar,\\
236~g de sal y 76~g de levadura respecto de las cantidades objetivo
};

% ---------- Efectos ----------
\coordinate (nivelEfectos) at ($(problema.north)+(0,2.1)$);

\node (efecto1) [efecto] at ($(nivelEfectos)+(-3.4,0)$) {
Menor disponibilidad\\del operario para\\otras actividades
};

\node (efecto2) [efecto] at ($(nivelEfectos)+(3.4,0)$) {
Masas con cantidades\\de ingredientes\\distintas de la receta
};

% ---------- Causas ----------
\node (causa2) [causa, below=1.8cm of problema] {
Ajuste final\\de la harina por\\encima del objetivo
};

\node (causa1) [causa, left=1.0cm of causa2] {
Cargas y traslados\\sucesivos desde\\el almacenamiento
};

\node (causa3) [causa, right=1.0cm of causa2] {
Azúcar, sal y\\levadura en una\\carga sin lectura
};

% ---------- Subcausas de la causa 1 ----------
\node (sub12) [sub, below=1.5cm of causa1] {
Tres o cuatro\\cargas de harina\\por ciclo
};

\node (sub11) [sub, left=0.3cm of sub12] {
Separación entre\\almacenamiento\\y pesaje
};

\node (sub13) [sub, right=0.3cm of sub12] {
Recorrido\\adicional para\\el azúcar
};

% ---------- Subcausas de la causa 3 ----------
\node (sub31) [sub, below=1.5cm of causa3, xshift=-1.45cm] {
Porción estimada\\por experiencia
};

\node (sub32) [sub, below=1.5cm of causa3, xshift=1.45cm] {
Envase común\\sin medida
};

% ---------- Flechas: problema -> efectos ----------
\coordinate (unionEf) at ($(problema.north)+(0,0.55)$);
\coordinate (unionEfIzq) at (efecto1.south |- unionEf);
\coordinate (unionEfDer) at (efecto2.south |- unionEf);

\draw[-] (problema.north) -- (unionEf);
\draw[-] (unionEfIzq) -- (unionEfDer);
\draw (unionEfIzq) -- (efecto1.south);
\draw (unionEfDer) -- (efecto2.south);

% ---------- Flechas: causas -> problema ----------
\coordinate (unionCausas) at ($(causa2.north)+(0,0.45)$);
\coordinate (unionCausasIzq) at (causa1.north |- unionCausas);
\coordinate (unionCausasCentro) at (causa2.north |- unionCausas);
\coordinate (unionCausasDer) at (causa3.north |- unionCausas);

\draw[-] (causa1.north) -- (unionCausasIzq);
\draw[-] (causa2.north) -- (unionCausasCentro);
\draw[-] (causa3.north) -- (unionCausasDer);
\draw[-] (unionCausasIzq) -- (unionCausasDer);
\draw (unionCausasCentro) -- (problema.south);

% ---------- Flechas: subcausas -> causa 1 ----------
\coordinate (unionSub1) at ($(sub12.north)+(0,0.55)$);
\coordinate (unionSub1Izq) at (sub11.north |- unionSub1);
\coordinate (unionSub1Der) at (sub13.north |- unionSub1);

\draw[-] (sub11.north) -- (unionSub1Izq);
\draw[-] (sub12.north) -- (unionSub1);
\draw[-] (sub13.north) -- (unionSub1Der);
\draw[-] (unionSub1Izq) -- (unionSub1Der);
\draw (unionSub1 -| causa1.south) -- (causa1.south);

% ---------- Flechas: subcausas -> causa 3 ----------
\coordinate (unionSub3) at ($(sub31.north)+(0,0.55)$);
\coordinate (unionSub3Izq) at (sub31.north |- unionSub3);
\coordinate (unionSub3Der) at (sub32.north |- unionSub3);

\draw[-] (sub31.north) -- (unionSub3Izq);
\draw[-] (sub32.north) -- (unionSub3Der);
\draw[-] (unionSub3Izq) -- (unionSub3Der);
\draw (unionSub3 -| causa3.south) -- (causa3.south);

% ---------- Coordenadas auxiliares ----------
\coordinate (leftEdge) at ($(sub11.west)+(-0.35,0)$);
\coordinate (rightEdge) at ($(sub32.east)+(0.35,0)$);
\coordinate (labelX) at ($(leftEdge)+(-0.45,0)$);
\coordinate (causaMid) at ($(causa1)!0.5!(sub12)$);

\coordinate (midEfProb) at ($(efecto1.south)!0.5!(problema.north)+(0,0.25)$);
\coordinate (midProbCausa) at ($(problema.south)!0.5!(causa2.north)+(0,0.25)$);

\path let \p1=(leftEdge),  \p2=(midEfProb)    in coordinate (l1L) at (\x1,\y2);
\path let \p1=(rightEdge), \p2=(midEfProb)    in coordinate (l1R) at (\x1,\y2);
\path let \p1=(leftEdge),  \p2=(midProbCausa) in coordinate (l2L) at (\x1,\y2);
\path let \p1=(rightEdge), \p2=(midProbCausa) in coordinate (l2R) at (\x1,\y2);

\path let \p1=(labelX), \p2=(efecto1)  in coordinate (lblEfecto)   at (\x1,\y2);
\path let \p1=(labelX), \p2=(problema) in coordinate (lblProblema) at (\x1,\y2);
\path let \p1=(labelX), \p2=(causaMid) in coordinate (lblCausa)    at (\x1,\y2);

% ---------- Líneas segmentadas ----------
\draw[dash pattern=on 6pt off 5pt, line width=0.6pt, -] (l1L) -- (l1R);
\draw[dash pattern=on 6pt off 5pt, line width=0.6pt, -] (l2L) -- (l2R);

% ---------- Etiquetas ----------
\node[banda] at (lblEfecto)   {Efectos};
\node[banda] at (lblProblema) {Problema\\central};
\node[banda] at (lblCausa)    {Causas};

\end{tikzpicture}
